\documentclass[10pt,a4paper]{amsart}
\usepackage[margin=19mm]{geometry}
\usepackage{amsmath,amssymb,mathtools}
\usepackage[scr=rsfs,cal=euler]{mathalfa}
\usepackage{xfrac}
\usepackage[unicode,hidelinks,bookmarks=false]{hyperref}
\newcommand{\ie}{i.e.\ }
\newcommand{\cf}{cf.\ }
\newcommand{\resp}{respectively\ }
\newcommand{\Subsection}{\subsection}
\providecommand{\nobreakdash}{\nobreak\hbox{-}}
\setlength{\emergencystretch}{2em}
\multlinegap0pt
\usepackage{bm}
\usepackage{bbm}
\usepackage{dsfont}
\usepackage{xspace}
\usepackage{cancel}
\usepackage{mathtools}
\usepackage{amsthm}
\makeatletter
\@ifundefined{lemma}{\newtheorem{lemma}{Lemma}}{}
\makeatother
\newcommand{\D}{\nabla}
\newcommand{\Ne}{\mathcal{N}_e}
\newcommand{\Ni}{\mathcal{N}_i}
\newcommand{\R}{\mathbb{R}}
\newcommand{\al}{\alpha}
\newcommand{\calf}{\mathcal{F}}
\newcommand{\dive}{\mathrm{div}}
\newcommand{\eps}{\mathbb{\varepsilon}}
\newcommand{\pa}{\partial}
\newcommand{\pt}{\partial_t}
\newcommand{\vp}{\varphi}
\title[Global Convergence and Error Estimates in Infinity-ion-mass ]{Global Convergence and Error Estimates in Infinity-ion-mass Limits for Bipolar Euler-Poisson System}
\author{Yachun Li, Shihao Wang, Liang Zhao}
\date{Source: arXiv:2211.16436v2 (CC BY 4.0)}
\begin{document}
\maketitle

\noindent\emph{Source and licence.} Global Convergence and Error Estimates in Infinity-ion-mass Limits for Bipolar Euler-Poisson System. Version arXiv:2211.16436v2 (CC BY 4.0), distributed under the Creative Commons Attribution 4.0 International licence, as recorded in the arXiv licence metadata for this exact version and shown on the abstract page https://arxiv.org/abs/2211.16436v2. Full rights notice: licenses/arxiv-2211.16436-CC-BY-4.0.txt.\par\smallskip

\noindent\textit{Editorial scope.} Counted unit: the Lemma whose source label is \texttt{14} in the article (source0.tex lines 956--961, source label \texttt{14}), delivered together with the article's own complete original proof of it (source0.tex lines 963--1129, one \texttt{proof} environment of 167 lines). No mathematics is written, repaired or re-derived: statement and proof are the article's own text line for line. Required context retained uncounted: none. Every definition, display and label of those lines is retained unchanged. Omitted from this package and not redistributed: 1--955, 962--962, 1130--1504. Nothing inside a retained excerpt is omitted or altered: every retained body is delivered exactly as the article's lines are written, so no omission receipt of any kind accompanies this package. Citations: the retained excerpts carry no citation command at all, so no bibliography entry of the article is transcribed and no reference is redistributed. Reference adaptation: none was needed and none was made; every reference inside the delivered unit resolves inside it, so no unresolved reference or citation is produced. Printed slips kept literal: no printed word, symbol, hypothesis or number of the counted statement or of its proof was altered. Layout and class: the article's own class is \texttt{amsart} with packages including bbm, amsfonts, graphicx, amssymb, mathrsfs, amsmath, amsthm, hyperref, soul, color, geometry, enumerate, xcolor, ulem; the wrapper is the lane's pinned \texttt{amsart} editorial wrapper at 10pt on A4, which restates the theorem-like environments the retained text needs and adds the article's own macro definitions that the retained text needs (\texttt{\textbackslash D}, \texttt{\textbackslash Ne}, \texttt{\textbackslash Ni}, \texttt{\textbackslash R}, \texttt{\textbackslash al}, \texttt{\textbackslash calf}, \texttt{\textbackslash dive}, \texttt{\textbackslash eps}, \texttt{\textbackslash pa}, \texttt{\textbackslash pt}, \texttt{\textbackslash vp}), with extra wrapper packages (bm, bbm, dsfont, xspace, cancel, mathtools). The delivered numbering is the unit's own. Assets: no retained span contains a figure environment, a graphics-inclusion command, a PDF-page inclusion, a table or a TikZ picture; no image file is copied into this package, the package declares no asset dependency and no image file is redistributed with it. Licence: version arXiv:2211.16436v2 of this article is distributed under the Creative Commons Attribution 4.0 International licence, as recorded in the arXiv licence metadata for this exact version and shown on the abstract page https://arxiv.org/abs/2211.16436v2. A full rights notice is delivered at licenses/arxiv-2211.16436-CC-BY-4.0.txt. Characters: every retained excerpt is pure ASCII, so the delivered engine, XeLaTeX with its default Latin Modern text font, reports no missing character.
\begin{lemma}\label{14} (Estimates for electron variables) It holds
\begin{equation}\label{regue}
    \|(\Ne,w_e,\mathcal{F})(t)\|_{s-1}^2 + \int_{0}^{t}\left\|w_e(\tau)\right\|^2_{s-1} d\tau 
    \leq  C\eps^2+C\delta\int_0^t\mathcal{D}(\tau)d\tau.
\end{equation}
\end{lemma}

\begin{proof} The error system for the electrons are of the form
\begin{equation*}
\begin{cases}
 \pt \mathcal{N}_e +\rho_e \dive w_e+\nabla\mathcal{N}_e\cdot u_e+\nabla\bar{\rho}_e\cdot w_e=- \mathcal{N}_e \dive\bar{u}_e,\\
 \pt w_e+(u_e \cdot \nabla) w_e+h'_e(\rho_e)\nabla\mathcal{N}_e+\mathcal{N}_e\nabla h'_e(\bar{\rho}_e)=-(w_e \cdot \nabla) \bar{u}_e
     -\mathcal{F}-w_e-r(\bar{\rho}_e,\mathcal{N}_e),
\end{cases}    
\end{equation*}
where
\begin{equation*}
    r(\bar{\rho}_e,\mathcal{N}_e)=(h'_e(\rho_e)-h'_e(\bar{\rho}_e)-h''_e(\bar{\rho}_e)\mathcal{N}_e)\nabla \bar{\rho}_e.
\end{equation*}

It can be written into the following first-order quasi-linear system:
\begin{equation}\label{16}
    \pt \mathcal{V}+\sum_{j=1}^3 A_j(\rho_e,u_e) \partial_{x_j} \mathcal{V}+ L(\bar{\rho}_e)\mathcal{V}=f,
\end{equation}
with 
\begin{equation*}
    \mathcal{V}=\begin{pmatrix}
    \mathcal{N}_e \\
    w_e
    \end{pmatrix}, \ \ \
    A_j(\rho_e,u_e)=\begin{pmatrix}
    u_{e}^j & \rho_e \xi_j^\top\\
    h'_e(\rho_e)\xi_j& u_{e}^j {\bf{I}}_3
    \end{pmatrix},\ \ \ 
    j=1,2,3,
\end{equation*}
and
\begin{equation*}
    L(\bar{\rho}_e)=\begin{pmatrix}
    0 & (\nabla \bar{\rho}_e)^\top\\
    \nabla h'_e(\bar{\rho}_e) & 0
    \end{pmatrix}, \quad 
    f=\begin{pmatrix}
    - \mathcal{N}_e \dive\bar{u}_e\\
    -(w_e \cdot \nabla) \bar{u}_e
     -\mathcal{F}-w_e-r(\bar{\rho}_e,\mathcal{N}_e)
     \end{pmatrix},
\end{equation*}
where $u_e=(u_e^1, u_e^2, u_e^3)$, ${\bf{I}}_3$ is the $3\times 3$ unit matrix and $\{\xi_k\}_{k=1}^3$ is the canonical basis of $\R^3$. 
For all $|\alpha| \leq s-1$, applying $\pa_x^\al$ to \eqref{16}, we obtain
\begin{equation}\label{17}
    \pt \pa_x^\al \mathcal{V}+\sum_{j=1}^3 A_j(\rho_e,u_e) \partial_{x_j} \pa_x^\al \mathcal{V}+ L(\bar{\rho}_e)\pa_x^\al \mathcal{V}=\pa_x^\al f+g^\alpha,
\end{equation}
in which
\begin{equation*}
    g^\alpha=\sum_{j=1}^3 A_j(\rho_e,u_e) \partial_{x_j} \pa_x^\al \mathcal{V}-\pa_x^\al \left ( \sum_{j=1}^3 A_j(\rho_e,u_e) \partial_{x_j} \mathcal{V} \right )+L(\bar{\rho}_e)\pa_x^\al \mathcal{V}-\pa_x^\al \left (L(\bar{\rho}_e)\mathcal{V} \right ).
\end{equation*}
Defining the symmetrizer $A_0(\rho_e)$ as follows
\begin{equation*}
    A_0(\rho_e)=\begin{pmatrix}
    h'_e(\rho_e) & 0\\
    0 & \rho_e {\bf{I}}_3
    \end{pmatrix},
\end{equation*}
and taking the inner product of \eqref{17} with $2A_0(\rho_e)\pa_x^\al \mathcal{V}$ in $L^2$ yields the following energy equality:
\begin{equation}\label{18}
\begin{aligned}
    \frac{d}{dt} \left< A_0(\rho_e)\pa_x^\al \mathcal{V}, \pa_x^\al \mathcal{V} \right> = & \left< \pt A_0(\rho_e)\pa_x^\al \mathcal{V}, \pa_x^\al \mathcal{V} \right>+\left<  B(\mathcal{V},\nabla\mathcal{V})\pa_x^\al \mathcal{V}, \pa_x^\al \mathcal{V} \right>\\
    & +2\left<  A_0(\rho_e)g^\alpha, \pa_x^\al \mathcal{V} \right>+2\left<  A_0(\rho_e)\pa_x^\al f, \pa_x^\al \mathcal{V} \right>\\
    \overset{def}{=}& J_1^\alpha+J_2^\alpha+J_3^\alpha+J_4^\alpha,
\end{aligned}
\end{equation}
with the natural correspondence of $J_1^\alpha$ to $J_4^\alpha$ and the matrix $B(\mathcal{V},\nabla\mathcal{V})$ is defined as follows

\begin{equation*}
\begin{aligned}
    B(\mathcal{V},\nabla\mathcal{V})&=\begin{pmatrix}
    \dive(h'_e(\rho_e)u_e) & (\nabla p'(\rho_e)-2h'_e(\rho_e)\nabla\bar{\rho}_e)^\top\\
    \nabla p'(\rho_e)-2\rho_e \nabla h'_e(\bar{\rho}_e) & \dive(\rho_e u_e){\bf{I}}_3
    \end{pmatrix}\\
    &\overset{def}{=} \begin{pmatrix}
    B_{11} & B_{12} \\
    B_{21}& B_{22}
    \end{pmatrix}.
\end{aligned}
\end{equation*}
For $J_1^\alpha$, we have
\begin{equation}\label{19}
| J_1^\alpha| \leq \left\|\pt A_0(\rho_e) \right\|_\infty \left\|\pa_x^\al \mathcal{V} \right\|^2 \leq C\delta \mathcal{D}(t).
\end{equation}
For $J_2^\alpha$, direct calculations give
\begin{equation*}
    J_2^\alpha=\left< B_{11} \pa_x^\al\mathcal{N}_e ,\pa_x^\al \mathcal{N}_e \right>+\left< B_{22} \pa_x^\al w_e ,\pa_x^\al w_e \right>+\left< (B_{12}+B_{21}^\top) \pa_x^\al\mathcal{N}_e ,\pa_x^\al w_e \right>,
\end{equation*}
where
\begin{equation}\label{21}
    \left< B_{11} \pa_x^\al\mathcal{N}_e ,\pa_x^\al \mathcal{N}_e \right> \leq 
    \left\| B_{11} \right\|_\infty \left\| \pa_x^\al \mathcal{N}_e \right\|^2 \leq C\delta \left\| \pa_x^\al \mathcal{N}_e \right\|^2,
\end{equation}
and
\begin{equation}\label{22}
    \left< B_{22} \pa_x^\al w_e ,\pa_x^\al w_e \right> \leq 
    \left\| B_{22} \right\|_\infty \left\| \pa_x^\al w_e \right\|^2 \leq C\delta \left\| \pa_x^\al w_e \right\|^2.
\end{equation}

From Taylor Expansion we have,
\begin{equation*}
    B_{12}(\rho_e,\nabla\rho_e)=B_{12}(\bar{\rho}_e,\nabla \bar{\rho}_e)+\frac{\partial B_{12}(\rho_e^a,\nabla \rho_e^b)}{\partial \rho_e}\mathcal{N}_e+\frac{\partial B_{12}(\rho_e^a,\nabla \rho_e^b)}{\partial (\nabla \rho_e)}\nabla \mathcal{N}_e,
\end{equation*}
and
\begin{equation*}
    B_{21}(\rho_e,\nabla\rho_e)=B_{21}(\bar{\rho}_e,\nabla \bar{\rho}_e)+\frac{\partial B_{21}(\rho_e^c,\nabla \rho_e^d)}{\partial \rho_e}\mathcal{N}_e+\frac{\partial B_{21}(\rho_e^c,\nabla \rho_e^d)}{\partial (\nabla \rho_e)}\nabla \mathcal{N}_e,
\end{equation*}
with $\rho_e^a$, $\rho_e^b$, $\rho_e^c$ and $\rho_e^d$ being between $\rho_e$ and $\bar{\rho}_e$.
Note that,
\begin{equation*}
\frac{\partial B_{12}}{\partial \rho_e}=\left( p_e'''(\rho_e)\nabla \rho_e-2h''_e(\rho_e) \nabla \bar{\rho}_e \right )^\top,\quad 
\frac{\partial B_{12}}{\partial (\nabla \rho_e)}=p_e''(\rho_e){\bf{I}}_3,
\end{equation*}
and
\begin{equation*}
\frac{\partial B_{21}}{\partial \rho_e}= p_e'''(\rho_e)\nabla \rho_e-2\nabla h'_e(\rho_e),\quad 
\frac{\partial B_{21}}{\partial (\nabla \rho_e)}=p_e''(\rho_e){\bf{I}}_3.
\end{equation*}
Also since
\begin{equation*}
    \left (B_{12}(\bar{\rho}_e,\nabla \bar{\rho}_e) \right)^\top=-B_{21}(\bar{\rho}_e,\nabla \bar{\rho}_e),
\end{equation*}
we have
\begin{equation}\label{20}
    \left< (B_{12}+B_{21}^\top) \pa_x^\al\mathcal{N}_e ,\pa_x^\al w_e \right> \leq \left\|B_{12}+B_{21}^\top \right\|_\infty \left\|\pa_x^\al \mathcal{N}_e\right\| \left\|\pa_x^\al w_e\right\| \leq C\delta \left (\left\|\pa_x^\al \mathcal{N}_e\right\|^2+\left\|\pa_x^\al w_e\right\|^2 \right).
\end{equation}
Combining \eqref{21}, \eqref{22} and \eqref{20} we conclude
\begin{equation}\label{23}
    |J_2^\alpha| \leq C\delta \mathcal{D}(t).
\end{equation}

For $J_3^\alpha$, by Moser-type inequality and Young's inequality we have
\begin{eqnarray}\label{24}
        &&\frac{1}{2}|J_3^\alpha| \nonumber \leq C\delta \mathcal{D}(t).
\end{eqnarray}

For $J_4^\alpha$, by Moser-type calculus inequalities,  we have
\begin{eqnarray}\label{31}
         J_4^\alpha \leq  -\left\langle 2\rho_e \pa_x^\al \calf , \pa_x^\al w_e \right\rangle-  \left\|\pa_x^\al w_e \right\|^2- C\delta\mathcal{D}(t). 
\end{eqnarray}
Noticing the equivalence of $\left< A_0(\rho_e)\pa_x^\al \mathcal{V}, \pa_x^\al \mathcal{V} \right>$ and $\left\|\pa_x^\al \mathcal{V}\right\|^2$ as $A_0$ is uniformly bounded, integrating \eqref{18} over $[0,t]$ and combining \eqref{19}, \eqref{23}, \eqref{24} and \eqref{31} yields
\begin{equation}\label{N_e}
    \left\|\pa_x^\al \mathcal{N}_e(t)\right\|^2+\left\|\pa_x^\al w_e(t)\right\|^2+ \int_{0}^{t}\left\|\pa_x^\al w_e(\tau)\right\|^2 d\tau+2\int_{0}^{t} \left\langle \rho_e \pa_x^\al \calf , \pa_x^\al w_e \right\rangle d\tau \leq  C\eps^2+C\delta \int_0^t \mathcal{D}(\tau)d\tau.
\end{equation}

Now it remains to estimate the last term on the left hand side of the above inequality. Since
\[
\left<\rho_e\pa_x^\al\calf,\pa_x^\al w_e\right>=\left<\pa_x^\al (\rho_e w_e), \pa_x^\al \calf\right>-\left<\pa_x^\al (\rho_e w_e)-\rho_e\pa_x^\al w_e,\pa_x^\al\calf\right>,
\]
in which 
\[
\left|\left<\pa_x^\al (\rho_e w_e)-\rho_e\pa_x^\al w_e,\pa_x^\al\calf\right>\right|\leq C\|\D \rho_e\|_{s-1}\|w_e\|_{s-1}\|\pa_x^\al \calf\|\leq C\delta \mathcal{D}(t).
\]
In addition, 
\begin{eqnarray*}
\left<\pa_x^\al (\rho_e w_e), \pa_x^\al \calf\right>=\left<\pt\pa_x^\al\Ne,\pa_x^\al(\vp-\bar{\vp})\right>-\left<\pa_x^\al(\Ne\bar{u}_e),\pa_x^\al \calf\right>,
\end{eqnarray*}
in which 
\[
\left|\left<\pa_x^\al(\Ne\bar{u}_e),\pa_x^\al \calf\right>\right|\leq C\|\bar{u}_e\|_s\|\Ne\|_{s-1}\|\pa_x^\al \calf\|\leq C\delta\mathcal{D}(t),
\]
and
\begin{eqnarray*}
\left<\pt\pa_x^\al\Ne,\pa_x^\al(\vp-\bar{\vp})\right>&=&\left<\pt\pa_x^\al(\Ne-\Ni),\pa_x^\al(\vp-\bar{\vp})\right>+\left<\pa_x^\al(\rho_iu_i),\pa_x^\al\calf\right>\nonumber\\
&\leq& -\dfrac{1}{2}\dfrac{d}{dt}\|\pa_x^\al\calf\|^2+C\|\eps^{-1}u_i\|_{s-1}^2+C\eps^2\|\calf\|_{s-1}^2.
\end{eqnarray*}
Substituting these estimates into \eqref{N_e} and summing the resulting equation for all $|\al|\leq s-1$ yield \eqref{regue}. 
\end{proof}

\end{document}
